(1) 芯内光线相对于轴的角度为 \(\theta_1\)，在芯—包层侧壁的入射角为 \(\pi/2-\theta_1\)。全反射要求 \(\sin(\pi/2-\theta_1)\ge n_2/n_1\)，即 \(n_1\sin\theta_1\le\sqrt{n_1^2-n_2^2}\)。端面Snell定律 \(n_0\sin\theta_0=n_1\sin\theta_1\)，故 \(n_0\sin\theta_{\max}=\sqrt{n_1^2-n_2^2}\)（若右侧不超过 \(n_0\)）。
(2) 数值孔径定义为 \(\mathrm{NA}=n_0\sin\theta_{\max}=\sqrt{n_1^2-n_2^2}\)；如果理论值大于 \(n_0\)，外侧所有可入射角均被几何接收，此时角域还受 \(\sin\theta_0\le1\) 限制。归一化频率 \(V=(2\pi a/\lambda)\sqrt{n_1^2-n_2^2}\)，其中 \(\lambda\) 是真空波长。
(3) 理想圆形弱导波阶跃光纤在 \(V<2.405\) 时通常只传播基模LP\(_{01}\)（考虑其偏振简并），高阶模式截止。增加芯半径或缩短波长会增大 \(V\) 并允许更多横向模式，因此高模式并非仅由折射率差决定。
(4) \(V\gg1\) 的弱导波多模估计 \(M\simeq V^2/2\)（按常见偏振计数约定），所以模式数近似随芯截面积及 \(\mathrm{NA}^2\) 增大。\(n_1\to n_2\) 时 \(\mathrm{NA}\to0\)、\(V\to0\)，横向束缚变弱；精确模结构由Maxwell矢量边值问题决定，几何接收角和大 \(V\) 模数式不适于所有极限。
